\providecommand{\Sp}{\mathrm{Sp}}
\providecommand{\Mp}{\mathrm{Mp}}
\providecommand{\Z}{\mathbb{Z}}
\providecommand{\Q}{\mathbb{Q}}
\providecommand{\R}{\mathbb{R}}
\providecommand{\C}{\mathbb{C}}
\providecommand{\HH}{\mathbb{H}}
\providecommand{\kBKM}{\kappa_{\mathrm{BKM}}}
\providecommand{\kch}{\kappa_{\mathrm{ch}}}
\providecommand{\kcat}{\kappa_{\mathrm{cat}}}
\providecommand{\kfib}{\kappa_{\mathrm{fiber}}}
\providecommand{\kHeis}{\kappa_{\mathrm{ch}}^{\mathrm{Heis}}}
\providecommand{\fg}{\mathfrak{g}}
\providecommand{\wt}{\mathrm{wt}}
\providecommand{\Lift}{\mathrm{Lift}}
\providecommand{\Dfive}{D_5}
\providecommand{\cO}{\mathcal{O}}

\section{Metaplectic boundary rows and Borcherds weights}
\label{sec:metaplectic-boundary-weights}

\subsection*{Metaplectic automorphy}

Let $\Gamma \subset \Sp_4(\Q)$ be a genus-two paramodular group and let
\[
  1 \longrightarrow \mu_2 \longrightarrow \widetilde{\Gamma}
  \longrightarrow \Gamma \longrightarrow 1
\]
be the pullback of Shimura's metaplectic cover.  A point of
$\widetilde{\Gamma}$ is represented by
$(\gamma,\varphi_\gamma)$ with
\[
  \gamma=\begin{pmatrix}A&B\\ C&D\end{pmatrix},\qquad
  \varphi_\gamma(Z)^2=\det(CZ+D),\qquad Z\in \HH_2 .
\]
A scalar-valued Siegel modular form of weight $r/2$ on this cover is a
holomorphic function $F\colon\HH_2\to\C$ satisfying
\[
  F(\gamma Z)=\nu(\widetilde{\gamma})\,
  \varphi_\gamma(Z)^r\,F(Z),
  \qquad \widetilde{\gamma}\in\widetilde{\Gamma},
\]
for a multiplier character $\nu\colon\widetilde{\Gamma}\to\C^\times$.
This is the precise meaning of half-integral weight: the form is a
section of a half-power automorphic line bundle on the cover.  Root
parity and a Borcherds--Kac--Moody denominator require separate
denominator data.

Gritsenko's additive lift applies after the Jacobi input and the Siegel
multiplier have been matched.  In this notation the lift has the form
\[
  \Lift_{\mathrm{Gr}}\colon
  J^{\mathrm w}_{k,t}(\rho)
  \longrightarrow
  M_k(\widetilde{\Gamma}^{(2)}_t,\nu),
\]
where $\rho$ is the Jacobi multiplier and $\nu$ is the induced Siegel
multiplier.  Cl\'ery--Gritsenko's diagonal-divisor theorem supplies
specific lifted products and their multiplier systems.  A denominator
algebra begins with a root lattice, a chamber, a Weyl vector, and signed
root multiplicities.

\subsection*{Three index systems}

Three finite lists meet at the level-one row.  They have different
indices and different normalisations.

\paragraph{Primitive CHL denominator ladder.}
The primitive Gritsenko--Nikulin denominator ladder is indexed by the
K3/elliptic compatible orders
\[
  N\in\{1,2,3,4,6\}.
\]
With the standard Borcherds input normalisation,
\[
\begin{array}{c|ccccc}
N&1&2&3&4&6\\ \hline
c_N(0)&10&8&6&4&2\\
\kBKM(\Phi_N^{\mathrm{BKM}})&5&4&3&2&1
\end{array}
\]
and $\kBKM(\Phi_N^{\mathrm{BKM}})=c_N(0)/2$.  The $N=1$ entry is
$\Delta_5$, of weight $5$.  The monic Gritsenko--Nikulin denominator is
\[
  \Dfive=64^{-1}\Delta_5,\qquad
  \operatorname{den}(\fg_{\Delta_5})=\Dfive(2Z).
\]
The scalar-square convention records
\[
  \Phi_{10}^{\mathrm{un}}=\Delta_5^2,\qquad
  \Phi_{10}^{\mathrm{OP}}=\Dfive^2=4096^{-1}\Delta_5^2 .
\]
Thus the primitive BKM denominator has weight $5$, while the scalar
square has weight $10$.

On the compact product row $Y=K3\times E$ the invariants come from
distinct constructions:
\[
\begin{array}{c|c|c}
\text{invariant}&\text{value}&\text{source}\\ \hline
\kcat(Y)&0&\chi(\cO_{K3})\chi(\cO_E)=2\cdot0\\
\kch(Y)&0&h^{0,\bullet}(Y)=(1,1,1,1)\\
\kHeis(Y)&3&\text{Heisenberg--Mukai specialisation}\\
\kBKM(\Delta_5)&5&c_1(0)/2=10/2\\
\kfib(Y)&24&\operatorname{rk}\widetilde H(K3,\Z)
\end{array}
\]
The product Hodge supertrace, the Heisenberg specialisation, the
Borcherds weight, and the Mukai fibre rank are four different
measurements.

\paragraph{Nikulin order-indexed twining.}
Nikulin's symplectic K3 orders are
\[
  N\in\{1,2,3,4,5,6,7,8\}.
\]
For the three orders outside the primitive CHL ladder the twined K3
Jacobi input has, in the order-indexed normalisation,
\[
\begin{array}{c|ccc}
N&5&7&8\\ \hline
\text{frame shape}&1^4 5^4&1^3 7^3&1^2 2^1 4^1 8^2\\
c_N(0)&4&2&2\\
\kBKM(\Phi_N)&2&1&1
\end{array}
\]
after the Borcherds input, cover, character, chamber, and product
normalisation have been fixed.  These are order-indexed K3 twining
constants.  The half-integral Cl\'ery--Gritsenko rows belong to the
diagonal-divisor catalogue below.

\paragraph{Gritsenko--Cl\'ery diagonal-divisor catalogue.}
The diagonal-divisor catalogue is indexed by paramodular triples
$(t,N_{\mathrm{dd}};k)$.  Its row data are
\[
\begin{array}{c|c|c|c}
\text{row}&(t,N_{\mathrm{dd}};k)&F_{t,N_{\mathrm{dd}}}
&c^{\mathrm{dd}}_{t,N_{\mathrm{dd}}}(0)\\ \hline
1&(1,1;5)&\Delta_5&10\\
2&(2,1;2)&F_{2,1}&4\\
3&(1,2;3)&F_{1,2}&6\\
4&(3,1;1)&F_{3,1}&2\\
5&(1,3;2)&F_{1,3}&4\\
6&(4,1;\tfrac12)&F_{4,1}&1\\
7&(1,4;\tfrac32)&F_{1,4}&3\\
8&(2,2;1)&F_{2,2}&2
\end{array}
\]
so the twice-weight tuple is
\[
  (10,4,6,2,4,1,3,2).
\]
The genuine half-integral rows are $(4,1;\tfrac12)$ and
$(1,4;\tfrac32)$, with half-integral multiplier-system data.
The only automatic overlap with the primitive CHL ladder is the
first row $\Delta_5$.  Any further overlap requires a row-map theorem
identifying the automorphic group, cover, character, divisor, Jacobi
input, and product normalisation on both sides.

\subsection*{Borcherds weight with fixed datum}

Let $f$ be a weakly holomorphic vector-valued modular form, or the
Jacobi form from which it is obtained, satisfying Borcherds' integrality
and principal-part hypotheses for a lattice of signature $(2,n)$.  Fix
the arithmetic group or metaplectic cover, character, Weyl chamber, and
normalisation of the leading product coefficient.  If $\Psi_f$ denotes
the resulting Borcherds product, then Borcherds' singular-theta theorem
gives
\[
  \wt(\Psi_f)=\frac{c_f(0)}{2}.
\]
Within this fixed datum define
\[
  \kBKM(\Psi_f):=\wt(\Psi_f)=\frac{c_f(0)}{2}.
\]
Changing the input form, squaring the product, replacing the cover, or
changing the character changes the datum.  Thus $\Delta_5$,
$\Delta_5^2$, a CHL product, an order-indexed K3 twining product, and a
Cl\'ery--Gritsenko diagonal-divisor row carry their own $c(0)$ data.
An equality between two such products is a theorem about the same
normalised automorphic product.

For half-integral Cl\'ery--Gritsenko rows the same formula is read on
the appropriate cover.  A value such as $1/2$ or $3/2$ is a modular
weight on that cover.  Cartan rank, simple-root multiplicities, and
bosonic/fermionic splitting are denominator-algebra data.

\subsection*{Denominator algebra criterion}

A Borcherds--Kac--Moody superalgebra attached to a boundary product
$F$ requires both a denominator datum and a multiplier.  The datum is:
\[
  (L,\mathcal C,\rho,P^{\mathrm{re}},W,\epsilon,\mu,F).
\]
Here $L$ is the Lorentzian product lattice, $\mathcal C$ a Weyl chamber,
$\rho$ the Weyl vector, $P^{\mathrm{re}}$ the set of real simple roots,
$W$ the reflection group, $\epsilon$ the reflection character, and
$\mu(\lambda)$ the integral signed root-multiplicity function.  The
denominator identity is the equality
\[
  F(Z)=e^\rho
  \prod_{\lambda\in L\cap\mathcal C}
  (1-e^\lambda)^{\mu(\lambda)}
  =
  \sum_{w\in W}\epsilon(w)\left(
    e^{w\rho}+\sum_a N(a)e^{w(\rho+a)}
  \right),
\]
with the same product-degree map on both sides.  Root parity is read
from a parity lift of the signed multiplicity function
\[
  \mu(\lambda)
  =\dim \fg_{\lambda,\bar 0}-\dim \fg_{\lambda,\bar 1}.
\]
When both even and odd parts are asserted, the lift to nonnegative
multiplicities is part of the structure.

A boundary product determines a Borcherds--Kac--Moody superalgebra after
the displayed datum, denominator identity, and parity lift are supplied.
A multiplier-system statement proves modularity of the product on its
cover.  A compact Hall interpretation then requires a Hall--Drinfeld
comparison with a source algebra.  On the K3 side this is the
Hall--Drinfeld branch attached to the Borcherds product.  The Mukai
self-mirror K3 Yangian is a separate stable-envelope branch.

\subsection*{Boundary geometry}

The fixed-origin elliptic automorphism construction used by the
$K3\times E$ CHL slice is available for
$N\in\{1,2,3,4,6\}$, equivalently for cyclotomic order
$\varphi(N)\mid 2$ on an elliptic curve.  Indeed the origin-fixing
automorphism group of a complex elliptic curve is the unit group of
its endomorphism ring: generically $\{\pm1\}$, at $j=1728$ the fourth
roots of unity, and at $j=0$ the sixth roots of unity.  Hence the
possible orders are exactly $\{1,2,3,4,6\}$.  Orders $5,7,8$ occur as
symplectic K3 automorphism orders and lie outside this elliptic
fixed-origin construction.

An $N$-torsion translation on an elliptic curve is a different datum.
The fixed-origin CM action used in the CHL specialisation must be
replaced by a new host theorem.  To use a translation quotient, or any
replacement host, one must prove the following statements:
\[
\begin{array}{ll}
\text{(i)}&\text{the quotient or stack has the required Calabi--Yau and
orientation data;}\\
\text{(ii)}&\text{the reduced stable-pairs or Donaldson--Thomas theory is
constructed on that host;}\\
\text{(iii)}&\text{the protected partition function equals the same
Borcherds product, with variables and chamber fixed;}\\
\text{(iv)}&\text{the one-particle Hall object or native }E_1\text{-chiral
}\Phi_3\text{ output has primitive index equal to the Jacobi input.}
\end{array}
\]
Before these four statements are proved, the boundary rows remain
automorphic products and, where a denominator datum is supplied,
Borcherds--Kac--Moody denominator products.  A reduced stable-pairs
partition function of a compact $K3\times E$-type CY$_3$ host is the
stronger geometric conclusion.

\subsection*{Source anchors}

Shimura 1975 constructs the metaplectic automorphy factor for
half-integral Siegel modular forms.  Howe 1979 identifies the Weil
representation governing the corresponding cover.  Borcherds 1998,
Theorem 13.3, gives the singular-theta weight formula
$\wt(\Psi_f)=c_f(0)/2$.  Gritsenko--Nikulin 1995/1998 give the
primitive denominator row containing $\Delta_5$.  Gritsenko 1999 fixes
the arithmetic-lift normalisation.  Gritsenko--Nikulin's product fixes
$\Dfive=64^{-1}\Delta_5$.  Cl\'ery--Gritsenko 2008/2011, Theorems
0.1, 1.2, and 3.1, give the diagonal-divisor genus-two catalogue, the
excluded ninth row, and the cusp-weighted product formula.  Mukai 1988
and Nikulin 1980 give the K3 symplectic automorphism orders.
Cheng--Duncan--Harvey 2014 and Eguchi--Hikami 2011 give the twined K3
moonshine data used in the order-indexed constants.
Oberdieck--Pandharipande and Oberdieck--Pixton fix the reduced
$K3\times E$ scalar square convention.
